% Work-MAT | 2. rocnik | pripravovana hodina 38
% Návaznost: uživatelské připravy.zip (hodiny 01–20), autorská učebnice na GitHubu.
% Realisticky.cz: didaktická inspirace v posloupnosti a typech úloh; zadání jsou vlastní.
\documentclass[11pt,a4paper]{article}
\input{../../sablony/pripravy/priprava.sty}
\nastavPripravu{2. ročník}{Logaritmické rovnice}{38}{Rovnice s logaritmickými vzorci}
\begin{document}
\blok{Motivace a návaznost}
Předvedeme sčítání a odčítání logaritmů uvnitř rovnice. Na tabuli pracuj v jednom sloupci: podmínky, spojení logaritmů, úprava argumentu, výsledek.
\poznamka{Orientační nabídka pro 45--55 minut: krátké zopakování a motivace, nový výklad, tři společné vzorové úlohy, samostatná práce. Vyber podle tempa třídy; zbývající zadání slouží jako rezerva.}
\blok{Výklad}
Vzorce o součinu, podílu a mocnině používáme jen pro kladné argumenty; nová rovnice s jediným logaritmem pak přechází na exponenciální tvar.
\blok{Společné příklady}
\begin{prikladReseniBox}[title={Příklad 1}]
Vyřeš, resp. urči: $\log_2(x-1)+\log_2(x+1)=3$.
\tcblower
\[\begin{aligned}x>1 \\ \log_2(x^2-1)=3 \\ x^2-1=8\;\Longrightarrow\;x=3\end{aligned}\]
\end{prikladReseniBox}
\begin{prikladReseniBox}[title={Příklad 2}]
Vyřeš, resp. urči: $\log_3(x-2)+\log_3 3=2$.
\tcblower
\[\begin{aligned}x>2 \\ 3(x-2)=9 \\ x=5\end{aligned}\]
\end{prikladReseniBox}
\poznamka{Nikdy neslučuj logaritmy s různými základy. Před úpravou ještě připomeň nesprávnou „větu“ o logaritmu součtu.}
\blok{Samostatné procvičování}
\begin{zadaniReseniBox}[title={Samostatná práce 1}]
\begin{enumerate}[label=\alph*)]
\item $\log_{2}(x-1)+\log_{2}2=1$.
\item $\log_{3}(x-2)+\log_{3}3=2$.
\item $\log_{5}(x-3)+\log_{5}5=3$.
\item $\log_{2}(x-4)+\log_{2}2=1$.
\item $\log_{3}(x-1)+\log_{3}3=2$.
\item $\log_{5}(x-2)+\log_{5}5=3$.
\end{enumerate}
\tcblower
\begin{enumerate}[label=\alph*)]
\item $x=2$.
\item $x=5$.
\item $x=28$.
\item $x=5$.
\item $x=4$.
\item $x=27$.
\end{enumerate}
\end{zadaniReseniBox}
\blok{Dobrovolné rozšíření}
\begin{zadaniReseniBox}[title={Rozšíření 1}]
\begin{enumerate}[label=\alph*)]
\item $Vytvoř vlastní úlohu tohoto typu a se spolužákem porovnej podmínky a výsledek.$.
\end{enumerate}
\tcblower
\begin{enumerate}[label=\alph*)]
\item Zkontroluj definice a argumenty..
\end{enumerate}
\end{zadaniReseniBox}
\blok{Úkol na doma}
U každého logaritmu vypiš podmínky argumentu a proveď zkoušku. U rovnic vyznač případně nepřípustné kořeny, u nerovnic odpověz intervalem.
\begin{zadaniReseniBox}[title={Domácí úkol}]
\begin{enumerate}[label=\alph*)]
\item $\log_{2}(x-3)+\log_{2}2=1$.
\item $\log_{3}(x-4)+\log_{3}3=2$.
\item $\log_{5}(x-1)+\log_{5}5=3$.
\item $\log_{2}(x-2)+\log_{2}2=1$.
\item $\log_{3}(x-3)+\log_{3}3=2$.
\item $\log_{5}(x-4)+\log_{5}5=3$.
\item $\log_{2}(x-1)+\log_{2}2=1$.
\item $\log_{3}(x-2)+\log_{3}3=2$.
\end{enumerate}
\tcblower
\begin{enumerate}[label=\alph*)]
\item $x=4$.
\item $x=7$.
\item $x=26$.
\item $x=3$.
\item $x=6$.
\item $x=29$.
\item $x=2$.
\item $x=5$.
\end{enumerate}
\end{zadaniReseniBox}
\poznamka{Domácí úkol obsahuje 8 položek; předpokládaná domácí práce přibližně 15--30 minut. Po větších blocích jsou vědomě promíchané dříve probrané typy.}
\end{document}
